\documentclass[11pt]{article}
\usepackage[margin=1in]{geometry}
\usepackage{amsmath, amssymb, amsthm}
\usepackage{booktabs}
\usepackage{enumitem}
\usepackage[hidelinks]{hyperref}

\newcommand{\cop}{\mathrm{cop}}
\newcommand{\pcf}{\mathrm{pcf}}
\newcommand{\code}[1]{\texttt{#1}}
\newcommand{\obs}{\mathrm{obs}}
\newcommand{\diag}{\mathrm{diag}}
\newcommand{\Var}{\mathrm{Var}}
\newcommand{\R}{\mathbb{R}}
\DeclareMathOperator*{\argmin}{arg\,min}

\title{Shared and object-specific factors\\
\large A hierarchical factor structure for surveys that observe different objects}
\author{Design note}
\date{September 18, 2026}

\begin{document}
\maketitle

\section{Purpose}\label{sec:purpose}

The estimated model has one factor space. Every survey's coefficients are
$\Lambda f_t$ plus measurement error, and whatever a survey shows that the space
cannot carry is called noise. Two symptoms of that, both established in the
HANK exercise, motivate this note:
\begin{enumerate}[leftmargin=*]
  \item The measurement-error multiplier on PSID's copula settles at $30$--$34$
        in every specification, i.e.\ the model treats the one survey that
        observes the full dependence structure as $30\times$ noisier than its
        bootstrap says. The conditional cells fail \emph{at} PSID's wave dates,
        not between them.
  \item The reconstructed series blip at PSID's wave quarters whenever the
        loadings are not PSID's own (every EM-based run; none of the
        complete-case runs): the arrival of $1{,}730$ coefficients every eighth
        quarter jolts the shared factors.
\end{enumerate}
Both are one conflict: PSID's wave-frequency information about the 3D copula
and the quarterly surveys' information about income want different factor
spaces, and the model has only one. The proposal is to give each \emph{object}
--- a marginal, a 2D copula slice, the 3D copula --- its own small block of
factors on top of the shared ones. Objects, not datasets, are the unit, because
an object has one set of dynamics and the datasets are partial, noisy views of
it: the income--wealth slice is seen by PSID, SCF and SIPP, and an
object-specific factor pools all three.

Section \ref{sec:objects} defines the objects and shows the partition is
already present in the coefficient indexing. Sections
\ref{sec:measurement}--\ref{sec:state} give the model. Section
\ref{sec:ident} is identification, Section \ref{sec:loadings} the estimation of
the new loadings, Section \ref{sec:recon} the reconstruction, Section
\ref{sec:why} why this does not repeat the block-diagonal failure, and Section
\ref{sec:minimal} the minimal version to run first.

\section{Objects and the partition of the coefficient vector}\label{sec:objects}

\paragraph{Measures and datasets.} There are $D$ measures (here $D = 3$:
consumption, income, wealth) and $J$ datasets (here $J = 5$). Dataset $j$
observes a subset $\mathcal{M}_j \subseteq \{1,\ldots,D\}$ of the measures on a
set of dates $\mathcal{T}_j \subseteq \{1,\ldots,T\}$.

\paragraph{Coefficients.} At each date, a dataset that observes all $D$
measures delivers
\begin{equation}
  \theta_{j,t} = \bigl(\theta^{\cop}_{j,t},\;\theta^{\pcf}_{j,t}\bigr) \in \R^{n},
  \qquad n = n_c + n_q ,
\end{equation}
the Legendre coefficients of the copula density and of the $D$ marginal
percentile functions. With $K$ Legendre degrees per dimension ($K = 12$,
degrees $0,\ldots,K-1$) and $G$ degrees per percentile function ($G = 12$):
\begin{itemize}[leftmargin=*]
  \item The copula tensor has $K^D = 1{,}728$ entries indexed by degree triples
        $(k_1,k_2,k_3)$, $k_d \in \{0,\ldots,K-1\}$. Entries with two or three
        zero indices are fixed by the uniform-marginal constraint (the
        ``immutable'' part): $1 + D(K-1) = 34$ of them. The free copula
        coefficients number $n_c = 1{,}728 - 34 = 1{,}694$.
  \item The marginals contribute $n_q = D\cdot G = 36$ coefficients.
\end{itemize}

\paragraph{The partition.} Classify each free copula coefficient by which of its
degree indices are zero. A zero index in dimension $d$ means the coefficient
does not vary with measure $d$, i.e.\ it belongs to the copula
\emph{marginalized over} $d$. Hence:
\begin{equation}
\begin{aligned}
  \mathcal{O}_{\mathrm{3D}}  &= \{(k_1,k_2,k_3): k_1,k_2,k_3 \ge 1\}, &&|\mathcal{O}_{\mathrm{3D}}| = (K-1)^3 = 1{,}331, \\
  \mathcal{O}_{\{d,d'\}}     &= \{(k_1,k_2,k_3): k_d,k_{d'} \ge 1,\ k_{d''} = 0\}, &&|\mathcal{O}_{\{d,d'\}}| = (K-1)^2 = 121 \ \text{each, three of them},\\
  \mathcal{O}_{d}^{\pcf}     &= \text{the $G$ percentile-function coefficients of measure $d$}, &&|\mathcal{O}_d^{\pcf}| = 12 \ \text{each, three of them}.
\end{aligned}
\end{equation}
Check: $1{,}331 + 3\cdot 121 + 3\cdot 12 = 1{,}730 = n$. The seven sets are
disjoint and exhaustive. Call them the \textbf{objects}, indexed by
$o \in \mathcal{O} = \{\mathrm{3D},\ \{C,I\},\ \{C,W\},\ \{I,W\},\ C^{\pcf},\ I^{\pcf},\ W^{\pcf}\}$,
with row-index sets $\mathcal{R}_o \subset \{1,\ldots,n\}$ and sizes
$n_o = |\mathcal{R}_o|$.

\paragraph{Which dataset observes which object.} Dataset $j$ observes object $o$
iff every measure the object involves is in $\mathcal{M}_j$:
\begin{equation}
  \obs(j) = \bigl\{ o : \mathrm{meas}(o) \subseteq \mathcal{M}_j \bigr\},
  \qquad
  \mathrm{meas}(\mathrm{3D}) = \{C,I,W\},\ \ \mathrm{meas}(\{d,d'\}) = \{d,d'\},\ \ \mathrm{meas}(d^{\pcf}) = \{d\}.
\end{equation}
For the HANK replicas:

\begin{center}
\begin{tabular}{llll}
\toprule
dataset & measures & objects observed & dates \\
\midrule
a (PSID)  & $C,I,W$ & all seven                                & 12 biennial waves \\
b (CPS)   & $I$     & $I^{\pcf}$                              & quarterly \\
c (CEX)   & $C,I$   & $C^{\pcf}, I^{\pcf}, \{C,I\}$          & quarterly \\
d (SCF)   & $I,W$   & $I^{\pcf}, W^{\pcf}, \{I,W\}$          & 9 triennial waves \\
e (SIPP)  & $I,W$   & $I^{\pcf}, W^{\pcf}, \{I,W\}$          & quarterly ($W$ annual) \\
\bottomrule
\end{tabular}
\end{center}

Two facts follow that drive everything below. The 3D object is observed by one
dataset on twelve dates. The $\{I,W\}$ slice --- the object the
\code{wealth\_by\_income} cells are computed from together with the 3D block ---
is observed by three datasets, one of them quarterly.

\section{Measurement equation}\label{sec:measurement}

\paragraph{Standardization.} As in the current model, each dataset's
coefficients are centred and scaled per object before estimation:
$\tilde\theta_{j,t} = S_j^{-1}(\theta_{j,t} - \mu_j)$, with $\mu_j \in \R^n$
the dataset's time mean and $S_j = \diag(s_{j,1},\ldots,s_{j,n})$ one scale per
(dataset, object), broadcast over the object's rows. Rows for objects
$\notin \obs(j)$ are missing (NaN) at every date.

\paragraph{States.} Let $f_t \in \R^{r}$ be the shared distributional factors
and, for each object $o$, $g_{o,t} \in \R^{r_o}$ its object-specific factors
($r_o \ge 0$; $r_o = 0$ means the object has none). Let $a_t \in \R^{q}$ be the
aggregate factors as now.

\paragraph{Loadings.} $\Lambda \in \R^{n\times r}$ is the shared loading matrix,
as now. For each object, $\Lambda_o \in \R^{n_o\times r_o}$ is its specific
loading matrix; embed it in $\R^{n\times r_o}$ as $\bar\Lambda_o$ with rows
$\mathcal{R}_o$ equal to $\Lambda_o$ and \emph{all other rows exactly zero}.

\paragraph{The equation.} For dataset $j$ at date $t$,
\begin{equation}
  \tilde\theta_{j,t}
  \;=\; H_{j,t}\Bigl( \Lambda f_t \;+\; \sum_{o\in\mathcal{O}} \bar\Lambda_o\, g_{o,t} \Bigr) \;+\; e_{j,t},
  \qquad e_{j,t}\sim N\bigl(0, \Sigma_j\bigr),
  \label{eq:meas}
\end{equation}
where $H_{j,t}$ is the diagonal selection matrix that keeps the rows observed by
$j$ at $t$ (zero rows for unobserved objects and for dates $t\notin\mathcal{T}_j$)
and $\Sigma_j = \diag(\sigma^2_{j,o}\mathbf{1}_{n_o})_{o}$ carries one variance
per (dataset, object) as now. Because $\bar\Lambda_o$ is zero outside
$\mathcal{R}_o$, the sum over $o$ contributes to dataset $j$'s rows only through
the objects $j$ observes: the $\{I,W\}$ factor enters PSID's, SCF's and SIPP's
income--wealth rows and nothing else; the 3D factor enters PSID's 3D rows only.

Stacking datasets, the observation vector $y_t \in \R^{N}$, $N = \sum_j n$,
has measurement matrix
\begin{equation}
  G_t = H_t\,\bigl[\; \Gamma\ \ \bar\Gamma_{o_1}\ \cdots\ \bar\Gamma_{o_7} \;\bigr],
  \qquad
  \Gamma = \begin{pmatrix}\Lambda\\ \vdots\\ \Lambda\end{pmatrix},\ \
  \bar\Gamma_o = \begin{pmatrix}\bar\Lambda_o\\ \vdots\\ \bar\Lambda_o\end{pmatrix}
  \quad(J\ \text{copies}),
\end{equation}
i.e.\ exactly the present $G_t$ with $\sum_o r_o$ extra columns. Under
\code{pre\_multiply} the same $\hat\Sigma^{-1/2}$ whitening is applied to
$y_t$ and to every column of $G_t$, new ones included. Annual ``tent''
observations (PSID flows) load $\tfrac14$ on each of four lags of $f$; the same
tent applies to $g_o$, which is why the specific states carry the same lag
chain as $f$ in Section \ref{sec:state}.

\section{State equation}\label{sec:state}

\paragraph{Shared block (unchanged).} With $F_t = (f_t, f_{t-1}, f_{t-2}, f_{t-3})$,
\begin{equation}
  \begin{pmatrix} f_t \\ a_t \end{pmatrix}
  = \begin{pmatrix} A & B \\ C & D \end{pmatrix}
    \begin{pmatrix} f_{t-1} \\ a_{t-1} \end{pmatrix} + \eta_t,
  \qquad \eta_t \sim N(0,\Omega),\quad
  \Omega = \begin{pmatrix}\Omega_f & \Omega_{fa}\\ \Omega_{af} & \Omega_a\end{pmatrix},
\end{equation}
$\Omega$ full (LKJ-correlated) as now; the stochastic-trend block, if used,
augments $f$ as before.

\paragraph{Object-specific blocks.} Each $g_{o,t}$ is an independent stationary
AR(1):
\begin{equation}
  g_{o,t} = \Phi_o\, g_{o,t-1} + \zeta_{o,t},
  \qquad \zeta_{o,t}\sim N\bigl(0,\ \diag(\omega_{o,1}^2,\ldots,\omega_{o,r_o}^2)\bigr),
  \qquad \Phi_o = \diag(\phi_{o,1},\ldots,\phi_{o,r_o}),\ \ |\phi_{o,\ell}| < 1,
  \label{eq:specific}
\end{equation}
with $\zeta_{o,t}$ independent across objects, across factors within an object,
and of $\eta_t$. In the full state vector
\begin{equation}
  x_t = \bigl(F_t,\ a_t,\ G_{o_1,t},\ \ldots,\ G_{o_7,t}\bigr),
  \qquad G_{o,t} = (g_{o,t}, g_{o,t-1}, g_{o,t-2}, g_{o,t-3}),
\end{equation}
the transition and innovation covariance are block-diagonal between the shared
block and every specific block:
\begin{equation}
  \mathcal{A} = \mathrm{blkdiag}\Bigl(\begin{pmatrix} A&B\\C&D\end{pmatrix}_{\text{with lag chain}},\ \Phi_{o_1}^{\text{lag}},\ \ldots\Bigr),
  \qquad
  \mathcal{Q} = \mathrm{blkdiag}\bigl(\Omega,\ \diag(\omega^2_{o_1}),\ \ldots\bigr).
\end{equation}
The specific blocks do not load on the aggregates ($C$-type rows zero) and are
not driven by them ($B$-type rows zero): by construction they carry what is
\emph{specific} to an object after the economy-wide movement has been taken by
$f$. Section \ref{sec:ident} explains why these zeros are not optional.

\paragraph{Dimensions.} With $r$ shared factors, $q$ aggregates and $\sum_o r_o$
specifics, the state has $4r + q + 4\sum_o r_o$ entries. Parameters added
relative to the current model: $2\sum_o r_o$ (one $\phi$ and one $\omega$ per
specific factor). Nothing is added to $\Omega$'s correlation block, which is
where extra \emph{shared} factors cost $O((r+q)^2)$.

\section{Identification}\label{sec:ident}

Three indeterminacies must be pinned down; the first is the one that decides
whether the extension means anything.

\paragraph{(i) Shared versus specific.} Without restrictions, an object's
specific factor is observationally equivalent to a shared factor whose loadings
happen to be zero outside $\mathcal{R}_o$; the model would be free to move
variation between $f$ and $g_o$, and a sampler facing a spare direction switches
it off (the 6-factor run did exactly that). The restriction that separates them
is the block-diagonality of \eqref{eq:specific}: $g_o$ is uncorrelated with
$\eta_t$ at all leads and lags and is not driven by $a_t$. Then the shared block
is the unique component that (a) is common to several objects and (b) comoves
with the aggregates, and $g_o$ is what remains in object $o$ after that
component is removed. Formally, for any two objects observed together in some
dataset, the cross-covariance of their coefficients at all leads and lags is
carried by $f$ alone,
\begin{equation}
  \mathrm{Cov}\bigl(\tilde\theta_{\mathcal{R}_o,t},\ \tilde\theta_{\mathcal{R}_{o'},t-h}\bigr)
  = \Lambda_{\mathcal{R}_o}\,\Gamma_f(h)\,\Lambda_{\mathcal{R}_{o'}}^{\top}
  \qquad (o\ne o'),
\end{equation}
so cross-object comovement identifies $\Lambda$ and the shared dynamics, and
the within-object autocovariance in excess of that identifies $\Lambda_o$,
$\Phi_o$, $\omega_o$.

\paragraph{(ii) Rotation and scale.} Within the shared block, rotation is fixed
by the PCA normalization of $\Lambda$ as now (orthogonal columns, ordered).
Within each specific block, the columns of $\Lambda_o$ are taken orthogonal and
ordered, and the scale is fixed by unit-variance scores on the columns the
loadings are estimated from --- the same convention as for $\Lambda$ (Section
\ref{sec:loadings}). With that, $\omega_{o,\ell}$ is the innovation sd of a
unit-scale factor and comparable across objects.

\paragraph{(iii) The 3D block is observed twelve times.} Between PSID waves
$g_{\mathrm{3D}}$ is unobserved. The likelihood therefore depends on
$\phi_{\mathrm{3D}}$ only through the eight-quarter transitions between
consecutive waves, i.e.\ through $\phi_{\mathrm{3D}}^{8}$ and the accumulated
variance $\omega_{\mathrm{3D}}^2(1-\phi^{16}_{\mathrm{3D}})/(1-\phi^2_{\mathrm{3D}})$,
from eleven wave-to-wave pairs. Persistence and innovation scale are nearly
unidentified separately. They must be fixed or given a tight prior:
\begin{equation}
  \phi_{\mathrm{3D}} \equiv \bar\phi \in [0.95, 0.99]
  \quad\text{(fixed, or } \phi_{\mathrm{3D}} \sim \text{a prior concentrated there)},
  \qquad \omega_{\mathrm{3D}} \text{ estimated}.
\end{equation}
This is the intended behaviour, not a limitation: a near-random-walk
$g_{\mathrm{3D}}$ \emph{holds} what PSID showed at the last wave until the next
one, blending it into the path instead of letting it decay in two quarters and
be discounted. The $\{I,W\}$ factor is observed quarterly by SIPP and its
$\phi$ can be estimated.

\paragraph{(iv) What is \emph{not} identified, and is not needed.} The
decomposition of a \emph{single} object's variance into shared and specific is
identified only through the shared block's comovement with other objects. An
object observed by one dataset that shares no measure with any other would have
an unidentified split; none of the seven objects is in that position (every
object shares at least one measure with a quarterly survey).

\section{Estimating the loadings}\label{sec:loadings}

\paragraph{Shared loadings.} $\Lambda$ is estimated exactly as now: complete-case
PCA, EM completion, or the Kalman-EM refinement on the rows other surveys
observe. Nothing in this note changes that step.

\paragraph{Specific loadings.} Let $\hat F$ be the shared factor path used to
estimate $\Lambda$ (the PCA scores, or the smoothed path under Kalman-EM), and
let $X_o$ be the pooled data matrix restricted to the rows $\mathcal{R}_o$ and to
the columns (dataset, date) at which object $o$ is observed. The residual after
the shared factors,
\begin{equation}
  R_o = X_o - \Lambda_{\mathcal{R}_o}\,\hat F_{\cdot,\mathrm{cols}(o)},
\end{equation}
is what $g_o$ has to explain. Take its leading $r_o$ principal components:
\begin{equation}
  \Lambda_o = U_o\,\mathrm{diag}(\lambda_{o,1},\ldots,\lambda_{o,r_o}),
  \qquad R_o \approx U_o\,\mathrm{diag}(\lambda_o)\,V_o^{\top},
\end{equation}
with the columns rescaled so the scores $V_o^{\top}$ have unit variance on the
columns used. Two remarks. First, $R_o$ pools every dataset observing $o$: for
$\{I,W\}$ that is PSID's twelve waves, SCF's nine and SIPP's quarters ---
several hundred columns, dominated by the surveys that see the object best.
Second, for the 3D block $R_{\mathrm{3D}}$ has PSID's twelve columns only, so
$\Lambda_{\mathrm{3D}}$ is a twelve-column PCA of PSID's residual --- the same
information as the complete-case loadings, which is unavoidable and is
precisely why the 3D factor must be persistent: it cannot be re-estimated
between waves, only carried.

\paragraph{Orthogonality in the sample.} Because $R_o$ is a residual on the
shared fit, its scores are orthogonal to $\hat F$ over the columns used. This is
the sample counterpart of restriction (i) and gives the sampler a starting
point consistent with the block-diagonal state equation.

\paragraph{Choosing $r_o$.} The rank cross-validation of
\code{doc/em\_rank\_selection.tex} applies object by object: hold out a share
of the observed cells of $R_o$, complete at each rank, score the held-out
cells. Objects whose curve is flat from $r_o = 1$ get none. Prior expectation
from the HANK runs: $r_{\mathrm{3D}} = 1$, $r_{\{I,W\}} = 1$, marginals $0$.

\section{Reconstruction}\label{sec:recon}

Given the smoothed states $\hat f_t$, $\hat g_{o,t}$ (and $\hat\tau_t$ if the
trend block is on), dataset $j$'s coefficient path is
\begin{equation}
  \hat\theta_{j,t} = \mu_j + S_j\Bigl( \Lambda(\hat f_t + \hat\tau_t) + \sum_{o\in\obs(j)} \bar\Lambda_o\,\hat g_{o,t} \Bigr) + \mathrm{trend}_{j,t},
\end{equation}
for every $t$, including dates $j$ did not observe. The cross-conditional cells
are then computed from $\hat\theta^{\cop}_{j,t}$ by the exact integration of
\code{compute\_cross\_conditional\_means\_continuous}, unchanged. Note what
enters \code{wealth\_by\_income}: the copula marginalized over consumption,
i.e.\ the $\{I,W\}$ slice, whose reconstruction is
$\Lambda_{\{I,W\}}\hat f_t + \Lambda_{\{I,W\}}^{\text{spec}}\hat g_{\{I,W\},t}$
--- the shared part plus a state informed by SCF and SIPP. The 3D factor enters
through the marginalization only indirectly (the slice is the degree-zero
consumption slice of the tensor; the 3D block's coefficients do not enter it).
So for the target cells the $\{I,W\}$ factor is the direct lever, and the 3D
factor's role is to stop PSID's 3D information from distorting $f$ --- which is
how it helps the cells and removes the blips at the same time.

\section{Why this is not the block-diagonal failure again}\label{sec:why}

The 4$\times$4 block-diagonal run gave the marginals and the copula separate
factors \emph{and nothing shared}. The marginal factors then had no route to
comove with the copula except through $\Omega$'s cross-block correlations,
which the LKJ prior shrinks toward zero; the filter's information about the
marginal states came from 36 rows instead of 1,730; three of four marginal
factors were estimated at zero innovation scale and the trend states absorbed
the movement.

Here the shared block stays. In the precision decomposition of
\code{doc/em\_and\_copula\_factors.tex}~\S5,
\begin{equation}
  \bigl[G_t^{\top}\Sigma^{-1}G_t\bigr]_{ff} = \sum_{i\ \text{observed}} \frac{\Lambda_i\Lambda_i^{\top}}{\sigma^2_{i}},
\end{equation}
every observed row still informs the shared factors, exactly as now; the
specific blocks \emph{add} states, they do not remove terms. An object's
specific factor being estimated at zero is then the correct outcome for an
object whose dynamics are fully shared --- shrinkage, not collapse --- and the
marginals are where that is expected.

\section{Priors}\label{sec:priors}

\begin{itemize}[leftmargin=*]
  \item $\phi_{o,\ell}$: for objects observed quarterly by some dataset, the
        same persistence prior family as the shared factors
        (\code{persistence\_prior = :per\_factor}); for the 3D block, fixed at
        $\bar\phi$ or $\mathrm{Beta}$ concentrated on $[0.95,0.99]$.
  \item $\omega_{o,\ell}$: half-Cauchy or the shared shock prior, with an
        optional global--local layer across specifics,
        $\omega_{o,\ell} \sim N^{+}(0,(\lambda_g\psi_{o,\ell})^2)$, so that
        specifics an object does not need switch off --- the same device as the
        trend block.
  \item Stationarity guard: the joint radius check extends to $\mathcal{A}$;
        since the specific blocks are diagonal AR(1)s with $|\phi|<1$, they
        cannot violate it, so the existing cap on the shared block suffices.
\end{itemize}

\section{The minimal version, and what to measure}\label{sec:minimal}

Start with two specific factors:
\begin{equation}
  r_{\{I,W\}} = 1,\qquad r_{\mathrm{3D}} = 1\ (\phi_{\mathrm{3D}} = 0.98\ \text{fixed}),\qquad
  r_o = 0 \ \text{otherwise},
\end{equation}
on top of the current $r = 5$ shared factors. State dimension $4\cdot 5 + 8 + 4\cdot 2 = 36$
(from 28); parameters $+3$ ($\phi_{\{I,W\}}$, two $\omega$'s). Targets, in
order of diagnostic value:
\begin{enumerate}[leftmargin=*]
  \item $\sigma^2_{a,\cop}$, the PSID copula multiplier: from $\approx 32$
        toward the prior mean. This is the direct test of the mechanism.
  \item The wave-quarter kink statistic on the smoothed factors
        (one-quarter deviation from the local trend at PSID wave quarters
        versus elsewhere, currently $1.4\times$ overall and $2.4\times$ on the
        first factor): toward $1\times$.
  \item The exact test at observation dates
        (\code{doc}: replica cells versus model cells at each survey's own
        dates): SCF and SIPP from $\approx 50$ toward their ceilings of $90+$.
  \item The \code{wealth\_by\_income} cells over all quarters, currently
        $41/33/43$ averaged over a, d, e.
  \item Guardrail: marginals and \code{income\_by\_wealth} must not fall.
\end{enumerate}
If the minimal version moves (1)--(3), the full seven-object version is
justified and its counts should be set by the per-object rank CV of Section
\ref{sec:loadings}; if it does not, the full version would not either.

\paragraph{Implementation note.} The stochastic-trend augmentation
(\code{StochasticTrend.jl}: \code{st\_augment\_G}, the collapsed smoother) is a
working template for appending states with their own dynamics and extending
$G_t$ column-wise; the object blocks differ from $\tau$ only in loading through
$\bar\Lambda_o$ instead of $\Lambda$ and in being AR(1) rather than
near-unit-root by default. The row sets $\mathcal{R}_o$ are recoverable from
the Legendre indexing already used by \code{retrieve\_cop\_and\_imm\_part} and
the marginalization in \code{compute\_cross\_conditional\_means\_continuous}.

\end{document}
